\section{System architecture}\label{sec:architecture}
\subsection{Separate physical, reporting and evaluation layers}
Figure~\ref{fig:architecture} depicts the redesigned system in separate lanes with orthogonal routes. The upper lane produces and releases observations; the middle lane forms immutable forecasts; the lower lane matches released feedback and updates calibration. An evaluator compartment remains outside these service lanes. This organisation distinguishes data movement from statistical operations and prevents arrows from passing through module labels. Table~\ref{tab:architecture-contract} specifies the inputs, outputs and information restrictions of each module, complementing the diagram rather than relying on visual proximity to imply an interface.

The physical source supplies speed, timestamp, sensor identity and original validity. In simulation it also supplies episode identity. An evaluator-side fault generator produces separate input and feedback release schedules before replay. Its permanent-loss markers and future schedule entries are hidden from the service. The input release queue updates a causal store, which constructs the seven-channel history described in Section~\ref{sec:data}. A separate feedback queue releases outcomes that can be matched to existing forecasts. The evaluator retains the uncensored truth and original mask exclusively for scoring.

\subsection{Frozen prediction and immutable issuance}
The forecasting lane contains a fitted backbone and noncrossing 0.05/0.50/0.95 head. For FAR-GW, gated dilated temporal convolutions interact with fixed and learned graph supports. The other backbone adapters replace this representation module while preserving the input interface and quantile output format. The median is the point forecast, and half the outer-quantile spread supplies a floored normalisation scale. Parameters remain frozen throughout calibration and test replay. Only the calibration state changes when eligible outcomes arrive.

Before publication of an interval, the issuance service creates a record containing sensor, issue time, target time, horizon, point, scale, context and interval bounds. The issued record is immutable. This ledger is necessary because a delayed target must be compared with the forecast actually made for it, not a later forecast generated after restoration. Model and calibrator versions and source manifests support auditing. Outcomes shared by several forecast horizons are legitimate, but duplicate ingestion of the same forecast--horizon residual is prohibited.

\subsection{Feedback matching, archives and calibration}
When a valid feedback packet arrives, the matcher retrieves its eligible forecast records and computes residuals from their stored centres and scales. The live archive retains target timestamps and issue contexts. A frozen archive from the calibration partition supplies fallback evidence. Target-time expiry is applied before querying live pools, so an old report cannot become recent merely because it was delivered in a backlog. No residual is inferred for a permanently missing report.

The calibration kernel queries a local graph group, applies target-age and context weights, measures timestamp-block support and computes a freshness gap. The real-data path uses capped local pools; the SUMO path also includes a global pool. These paths are explicit in the diagram and algorithm because their sampling and pooling estimands differ. The mixture shifts mass towards frozen evidence when live support is stale or insufficient. Signed-tail quantiles, bounded inflation and nonnegative clipping then produce the final interval. The output is stored before evaluator scoring can occur.

\subsection{Execution and reproducibility boundary}
Replay advances in five-minute bins: release packets, update causal inputs, ingest feedback, expire live records, query the frozen predictor and issue intervals. Evaluator-only scoring compares saved intervals with originally valid outcomes. Figure~\ref{fig:architecture} shows scoring as a terminal branch without a return edge to the service. This ordering is the core protection against information leakage, not a claim that every proposed weight is statistically optimal.

The implementation checks source manifests and chunk hashes before accepting completed evaluations. Resumed runners may revisit artifact manifests without discarding verified work. This is an execution property rather than an additional statistical replicate. Similarly, the gateway smoke tests establish limited replay-equivalence behaviour, not production latency, network reliability or operational readiness. Table~\ref{tab:architecture-contract} makes those boundaries reviewable: the source and evaluator may know truth, but the predictor and calibrator receive only the information already released to their respective channels.
